\section {Using the programs in executable form}
\typeout {Section 2 ***}

At the high-level end, {\tt LatMRG} provides programs in executable form.
These programs read their data in files or work interactively with the
user. Appendix A describes their use in more detail.
In this section, we give examples.
The user can also tailor his own programs using the lower-level
tools offered by the different modules of {\tt LatMRG}.
This is discussed in the next section.

The program {\tt maxper} checks whether a given generator has maximal period.
The program {\tt findmk} can find values of $m$ and $k$ such that
$r = (m^k-1)/(m-1)$ is prime. 
The programs {\tt latl} and {\tt lats} perform standard lattice and/or 
spectral tests for that generator.
The only difference between the latter two programs is that in {\tt latl}, 
the multipliers and basis components are implemented in {\tt LONGINT} and
{\tt LONGREAL}, respectively, while in {\tt lats}, they are implemented as
{\tt SuperInteger}'s.
More comprehensive (and larger) programs perform computer searchs
to seek the ``best'' generators of a given type, according to maximal period
and lattice structure criteria.
These programs also come in two versions.
In {\tt seekl}, the multipliers and basis components are implemented in
{\tt LONGINT} and {\tt LONGREAL}, respectively,
while in {\tt seeks}, both are implemented as {\tt SuperInteger}'s.

The program {\tt menumrg}  is an interactive program which allows the
user to build a lattice basis for a given generator, and then
manipulate, compare, or reduce lattice bases.

We now give a few concrete examples of data files and results.
The timings given are from runs on a SUN SparcStation 20, under SunOS 5.4,
using version 4.5.1 of the MCS Modula-2 compiler \cite{iJAG92a}.
These results were obtained with the `February 1996' version of {\tt LatMRG}.


%%%%%%%%%%%%%%%%%%%%%%%%%%%%%%%%%%%%%%%%%%%%%%%%%%%%%%%
%  This is in file  fish31.dat
%
\setbox0=\vbox {\hsize = 6.3in
\code\smallc

   FALSE                        ReadGenFile
   1                            J  (Nb. of components) 
   MRG                          Gener. type
   2147483647  1 0              m  (modulus) 
   1                            k  (order) 
   742938285                    a  (multiplier) 
   2  30                        MinDim   MaxDim
   BeyerSpectral  Hermite       LatticeInfo
   Full                         LatticeType
   1  1                         LaGroupSizes   Spacings
   NoVerify                     VerifyBB
   1000000                      MaxNodesBB
   TEX                          ResultForm

\endcode }

\begin{figure}[hbt] 
\centerline {\boxit{\box0}}
\caption { Example of a data file for the program {\tt latl},
           in file {\tt fish31.dat}. }
\label {lat1.dat}
\end{figure}

%%%%%%%%%%%%%%%%%%%%%%%%%%%%%%%%%%%%%%%%%%%%%%%%%%%%%%%

%  This is in file  fish31.tex
%
\begin {table}[hbt] 
\caption { Results of {\tt latl} in file {\tt fish31.tex}. }
\label {lat1.tex}
\medskip
\centering
\begin {tabular}{|rlllr|}
\hline
 &&&&  cumulative \\
   $t$ &      $d_t$   &      $q_t$  &      $S_t$  &  cpu (sec)\\
\hline
    2  &   2.31556E-5 &     0.84858 &     0.86725 &      0.02 \\
    3  &   8.02308E-4 &     0.90348 &     0.86068 &      0.03 \\
    4  &   4.52795E-3 &     0.88522 &     0.86270 &      0.04 \\
    5  &      0.01328 &     0.79661 &     0.83195 &      0.05 \\
    6  &      0.02586 &     0.85416 &     0.83415 &      0.06 \\
    7  &      0.05530 &     0.59317 &     0.62392 &      0.07 \\
    8  &      0.06820 &     0.64472 &     0.70666 &      0.09 \\
    9  &      0.10600 &     0.67859 &             &      0.10 \\
   10  &      0.10847 &     0.73126 &             &      0.13 \\
   11  &      0.16903 &     0.63121 &             &      0.15 \\
   12  &      0.24254 &     0.58925 &             &      0.19 \\
   13  &      0.24254 &     0.58080 &             &      0.22 \\
   14  &      0.24254 &     0.63681 &             &      0.26 \\
   15  &      0.24254 &     0.67276 &             &      0.31 \\
   16  &      0.24254 &     0.76358 &             &      0.39 \\
   17  &      0.24254 &     0.78289 &             &      0.48 \\
   18  &      0.25000 &     0.74582 &             &      0.64 \\
   19  &      0.26726 &     0.75484 &             &      1.14 \\
   20  &      0.26726 &     0.81215 &             &      1.71 \\
   21  &      0.26726 &     0.78067 &             &      3.15 \\
   22  &      0.28868 &     0.88552 &             &      5.77 \\
   23  &      0.28868 &     0.86355 &             &      8.37 \\
   24  &      0.30151 &     0.91232 &             &     13.46 \\
   25  &      0.30151 &     0.86634 &             &     26.85 \\
   26  &      0.30151 &     0.84667 &             &     48.06 \\
   27  &      0.30151 &     0.87238 &             &     80.57 \\
   28  &      0.30151 &     0.88873 &             &    126.43 \\
   29  &      0.31623 &     0.88771 &             &    200.71 \\
   30  &      0.31623 &     0.88201 &             &    307.96 \\
\hline
\end {tabular}
\end {table}

%%%%%%%%%%%%%%%%%%%%%%%%%%%%%%%%%%%%%%%%%%%%%%%%%%%%%%
\subsection {An example with the program {\tt latl}}

Figure \ref{lat1.dat} gives an example data file for the program {\tt latl}.
The corresponding results appear in Table~\ref{lat1.tex}.
To call the program and produce these results, type ``{\tt latl fish31}'',
assuming that the data are in file ``{\tt fish31.dat}''.
The results will be in file ``{\tt fish31.tex}'', which produces
Table~\ref{lat1.tex} after going through \LaTeX.
If {\tt TEX} was replaced by {\tt Terminal} in the data file,
the results would rather be displayed on the screen.
See the description of the program {\tt latl} for more details on how
to set up the data files.
Note that the first column in the data file gives the data values themselves,
while the second column contains comments describing the meaning of these
values.
In that example, the Beyer and spectral tests are applied to
the LCG of order 1 with $m = 2^{31}-1 = 2147483647$ and $a = 742938285$,
suggested by Fishman and Moore \cite {rFIS86a}.
The last column indicates the (cumulative) cpu time.
The total cpu time to compute all the Beyer quotients $q_t$ and
distances $d_t$ between hyperplanes in dimensions 2 to 30 was approximately
308 seconds, i.e., approximately 5 minutes.
Computing the Beyer quotients for this example turns out to be
much more expensive than computing only the distances between the hyperplanes.
For example, to compute only the $d_t$'s (and not the $q_t$'s)
in dimensions 2 to 30, it takes less than 3 seconds on the same machine!
We also computed $d_t$ and $q_t$ in dimensions 2 to 30 with formal
verification using error bounds for all floating point numbers
(see the module {\tt REDBAS2}).  This took approximately 15 minutes
and the results were the same (except for the CPU times)
as in Table~\ref{lat1.tex}.
These times are representative of what happens in general for examples 
of that size.
The program {\tt lats} can also be applied to the same data file and the
results will be the same, except for the cpu times.
It takes more time to run, but it can accept much larger moduli
and multipliers.


%%%%%%%%%%%%%%%%%%%%%%%%%%%%%%%
%  This is in file  seek15.dat
%
\setbox0=\vbox {\hsize = 6.0in
\code\smallc

 FALSE                           Read generators from file
 1                               J
  MRG                             Gener. type
  32749  1 0                      Modulo m
  2                               Order k
  TRUE                            Maximal period
  FALSE                           Implem. condition
  Decomp                          Factors of m-1
  DecompWrite seek15r.fac         Factors of r
  1                               b1
  180                             c1
  -180                            b2
  -1                              c2
  Exhaust                         Search method: exhaustive search
 1 3 8                           NbCat  Dim(0)  Dim(1)
 0.0  1.0                        Min and max merit values
 2                               Nb of retained generators
 Spectral  Hermite               Merit criterion
 Spectral                        Information to print
 Full                            Lattice type (analyzed)
 1  1                            Lacunary ind. group sizes,  spacing
 NoVerify                        Verify basis reduction with error bounds.
 1000000                         Max nb of nodes in each BB
 0.1                             Time limit: 0.1 hour
 12345 98765                     S1 S2 : seeds for random number generator
 RES                             Results File

\endcode }

\begin{figure}[hbt] 
\centerline {\boxit{\box0}}
\caption { Example of a data file for the program {\tt seekl},
           in file {\tt seek15.dat}. }
\label {seekl1.dat}
\end{figure}

%%%%%%%%%%%%%%%%%%%%%%%%%%%%%%%
%  This is in file  seek15.res
%
\setbox1=\vbox {\hsize = 5.0in
\code\smallc
 SEARCH FOR GOOD MRGs OF ORDER 2
-----------------------------------------

 DATA
 ----
   Modulus  m = 32749
   Order    k = 2
   n_j        = 1
   rho_j = (m_j)^k_j-1 = 1072497000
   Factors of m-1   :  Decomp 
   Factors of r     :  DecompWrite seek15r.fac
   Bounds : a1 from : 1
                 to : 180
            a2 from : -180
                 to : -1
   Search method                      : EXHAUST 
   Implem. cond.  a_i (m mod a_i) < m : NO
   Maximum period required            : YES

   Merit criterion                    : M_8
   Seeds for RNG                      : 12345,  98765
   Verify Branch-and-bound            : NoVerify
   Maximum nodes in branch-and-bound  : 1000000
   Lattice Type                       : Full

 RESULTS 
 -------
   Values of a2 tried                    :        180
   Values of a2 primitive element        :         66
   Polynomials with a2 primitive element :      11880
   Primitive polynomials                 :       4638
   Nb. of polynomials to examine         :      32400
   Nb. Generators conserved              :          2
   Total CPU time (after setup)          : 0:00:13.06
 -----------------------------------------------------
1. 
   a_1 =              180
   a_2 =             -176

    t         d_t             q_t           S_t
 -----------------------------------------------------
    3      3.97223E-3                     0.21911
    4      7.13449E-3                     0.65130
    5         0.03511                     0.36151
    6         0.03558                     0.68076
    7         0.09285                     0.41035
    8         0.09285                     0.56613

   Merit =    0.21911  = S_3

2. 
   a_1 =              180
   a_2 =             -175

    t         d_t            q_t            S_t
 -----------------------------------------------------
    3      3.98327E-3                     0.21850
    4      7.39990E-3                     0.62794
    5         0.02824                     0.44953
    6         0.03716                     0.65170
    7         0.05717                     0.66649
    8         0.07715                     0.68130

   Merit =    0.21850  = S_3

\endcode }

\begin{figure}[hbt] 
\centerline {\boxit{\box1}}
\caption { Results of program {\tt seekl} in file {\tt seek15.res}. }
\label {seekl1.res}
\end{figure}


%%%%%%%%%%%%%%%%%%%%%%%%%%%%%%%%%%%%%%%%%%%%%%%%%%%%%%%
\subsection {Examples with the programs {\tt seekl} and {\tt seeks}}

Figure \ref{seekl1.dat} gives an example of a data file for {\tt seekl}.
It will perform an exhaustive search among the 32400 generators of order
$k=2$ with modulus $m= 32749$, $1\le a_1\le 180$, and $-180\le a_2\le -1$.
The criterion is $M_8$.
The two best generators
will be retained, provided their values of $M_8$ are at least 0.2.
The values of $m-1$ and $r$ will be decomposed by the program.
The factors of $r$ will be written in file ``{\tt f2seek1.fac}'',
while those of $m-1$ will not be kept.
The results of that program appear in Figure~\ref{seekl1.res}
(this is an actual printout, for illustration).

Figure \ref{seeks2.dat} gives another data file example, this time
for {\tt seeks}.
It asks for a random search for good generators of order
$k=5$ with modulus $m= 2^{63} - 711$, $1\le a_1\le 2^{63}-712$,
$a_2 = a_3 = a_4 = 0$, $-2^{63}+712\le a_5\le -1$, for which
$a_i (m\mod a_i) < m$ for $i=1,5$, and which have maximal period.
The criterion is $S_8$.
Only the best generator will be retained.
The factorization of $m-1$ will be read in file ``{\tt seek63.fac}'' and
$r$ is prime.
For the random search, we will examine 1000 subregions of dimensions
$(4\times 1\times 1\times 1\times 4)$, that is, a total of 16000 generators
(4000 values of $a_5$), if time permits.
We give the program a cpu time-limit of 3 hours.
A partial view of the results file is given in Figure~\ref{seeks2.res}.
The program reached the time-limit of 3 hours before examining
all the generators it was asked for.
It then stopped and reported the best generator it had found so far.

Note that we are not recommending any of these particular generators.
These examples are only to illustrate the capabilities of LatMRG.

%%%%%%%%%%%%%%%%%%%%%%%%%%%%%%%%%%%%%%%%%%%%%%%%%%%%%%%
%  This is in file  seek63.dat
%
\setbox0=\vbox {\hsize = 6.1in
\code\smallc

     FALSE                          ReadGenFromFile
     1                              J
      MRG                            TypeGen
      2  63  -711                    m
      5                              k
      TRUE                           PerMax
      TRUE                           ImplemCond
      Read  seek63.fac               Factors of m-1
      Prime                          Factors of r
      1                              b_1
      2  63  -712                    c_1
      0                              b_2
      0                              c_2
      0                              b_3
      0                              c_3
      0                              b_4
      0                              c_4
      -2  63  -712                   b_5
      -1                             c_5
      Random 1000 4 4               Search: random, 1000 regions, H=4, Hk=4
     1 6 12                         C  Dim(0)  Dim(1)
     0.0  1.0                       Minimal and maximal merit values
     1                              Nb of retained generators
     Spectral       Hermite         Criterion
     BeyerSpectral  Hermite         Information to print
     Full                           Lattice type (analyzed)
     1  1                           Lacunary ind. group sizes,  spacing
     NoVerify                       Verify BB with error bounds
     1000000                        Max nb nodes in each BB
     3.0                            cpu time-limit (hours)
     12345 98765                    S1 S2 : seeds for random number generator
     RES                            Results File

\endcode }

\begin{figure}[hbt] 
\centerline {\boxit{\box0}}
\caption { Example data file for program {\tt seeks},
           in file {\tt seek63.dat}. }
\label {seeks2.dat}
\end{figure}

%%%%%%%%%%%%%%%%%%%%%%%%%%%%%%%%%%%%%%%%%%%%%%%%%%%%%%%
%  This is in file  seek63.res
%
\setbox1=\vbox {\hsize = 5.5in
\code\smallc

 SEARCH FOR GOOD MRGs OF ORDER 5
 ----------------------------------------

 \vdots

 RESULTS 
 -------
 ****  Program aborted: reached time limit.  **** 
     There will be no verification of results.
          time limit :       3.00 hour(s) 
   Values of a5 tried                    :        277
   Values of a5 primitive element        :        134
   Polynomials with a5 primitive element :        536
   Primitive polynomials                 :        103
   Nb. Generators conserved              :          1

   Total cpu time (after setup)          : 3:00:01.76
 -----------------------------------------------------

1. 
   a1 = 53663276718
   a2 = 0
   a3 = 0
   a4 = 0
   a5 = -5144244197

    t        d_t          q_t            S_t
 -----------------------------------------------------
    6     1.855E-11                   6.5589E-6
    7     1.855E-11                   1.1384E-3
    8     1.855E-11                     0.05347
    9     2.991E-11              
   10     3.360E-10              
   11     2.6017E-9              
   12     1.2892E-8              

   Merit = 6.55894E-6  = S_6

\endcode }

\begin{figure}[hbt] 
\centerline {\boxit{\box1}}
\caption { Results of program {\tt seeks} in file {\tt seek63.res}. }
\label {seeks2.res}
\end{figure}

\clearpage





